%% kapa: exam=1 round=31 year=2020 period=1 form=A subject=real_estate subject_ko=부동산학원론 seq=5 number=85
%% source: Q-Net 감정평가사 기출문제 (cst00309) artlSeq=5207219 file=extracted/1008/5207219/A형/A형/2020년도 제31회 감정평가사 1차 시험_1교시_A형.pdf
%% conversion: kordoc 4.24.51 -> tools/gen1.py (auto); review=proofread-sonnet; flags=none
\begin{question}{85}
감정평가 실무기준상 권리금 감정평가방법에 관한 설명으로 옳지 않은 것은?
\begin{choices}
  \item 권리금을 감정평가할 때에는 유형ㆍ무형의 재산마다 개별로 감정평가하는 것을 원칙으로 한다.
  \item 권리금을 개별로 감정평가하는 것이 곤란하거나 적절하지 아니한 경우에는 일괄하여 감정평가할 수 있으며, 이 경우 감정평가액은 유형재산가액과 무형재산가액으로 구분하지 않아야 한다.
  \item 유형재산을 감정평가할 때에는 주된 방법으로 원가법을 적용하여야 한다.
  \item 무형재산을 감정평가할 때에는 주된 방법으로 수익환원법을 적용하여야 한다.
  \item 유형재산과 무형재산을 일괄하여 감정평가할 때에는 주된 방법으로 수익환원법을 적용하여야 한다.
\end{choices}
\end{question}
\answer{2}
